\documentclass[conference]{IEEEtran}
\IEEEoverridecommandlockouts
\usepackage{cite}
\usepackage{amsmath,amssymb,amsfonts}
\usepackage{algorithmic}
\usepackage{graphicx}
\usepackage{textcomp}
\usepackage{xcolor}
\def\BibTeX{{\rm B\kern-.05em{\sc i\kern-.025em b}\kern-.08em
    T\kern-.1667em\lower.7ex\hbox{E}\kern-.125emX}}

\begin{document}

\title{MoodSync: Real-Time Facial Emotion Recognition\\with Adaptive Music Recommendation}

\author{
\IEEEauthorblockN{Salman Farsi}
\IEEEauthorblockA{\textit{Dept. of Computer Science and Engineering} \\
\textit{BRAC University}\\
Dhaka, Bangladesh \\
salman.farsi@g.bracu.ac.bd}
\and
\IEEEauthorblockN{Rifaat Nuha}
\IEEEauthorblockA{\textit{Dept. of Computer Science and Engineering} \\
\textit{BRAC University}\\
Dhaka, Bangladesh \\
rifaat.nuha@g.bracu.ac.bd}
}

\maketitle

% -------------------------------------------------------
\begin{abstract}
Emotion recognition from facial expressions plays a vital role in enabling
intelligent human computer interaction systems. Despite advances in deep
learning, achieving high classification accuracy on real world data while
maintaining practical deployability remains a significant challenge. This
study presents MoodSync, a real time facial emotion recognition system built
on a balanced subset of the FER-2013 dataset, comprising 10,000 training
images and 6,043 test images across five emotion categories: angry, happy,
neutral, sad, and surprise. Three convolutional neural network architectures
were evaluated and compared, namely a custom CNN, DenseNet121, and FERNET,
a deep VGG style network, to identify the most effective model for five
class emotion classification. All models were trained on preprocessed
48 by 48 pixel grayscale images with online data augmentation applied to
improve generalization. Performance was measured using accuracy, weighted
precision, weighted recall, weighted F1 score, and macro average AUC.
FERNET achieved the highest classification accuracy of 71.45\%, a weighted
F1 score of 71.60\%, and an AUC of 0.922, outperforming both the custom
CNN (64.80\%) and DenseNet121 (56.02\%). Beyond emotion classification,
MoodSync integrates a rule based music recommendation module that dynamically
suggests genre appropriate tracks aligned with the detected emotional state,
forming a complete end to end affective computing application suitable for
real time webcam deployment.
\end{abstract}

\begin{IEEEkeywords}
Facial Emotion Recognition, Convolutional Neural Network, FER-2013,
FERNET, VGGNet, Deep Learning, Affective Computing, Music Recommendation,
Human Computer Interaction
\end{IEEEkeywords}

% -------------------------------------------------------
\section{Introduction}

Human emotions are complex psychological states that manifest through facial
expressions, body language, and vocal cues. Of these channels, facial
expressions are the most universally consistent and visually accessible,
making them the primary target for automated affective computing systems.
The ability to automatically detect and interpret facial emotions in real
time has substantial implications across a wide range of applications,
including mental health monitoring, adaptive e-learning platforms, driver
drowsiness detection, and intelligent human computer interaction \cite{b1}.
As computing systems become increasingly embedded in daily life, enabling
machines to perceive and respond to human emotional states represents a
critical frontier in artificial intelligence research.

The FER-2013 dataset, one of the most widely used benchmarks for facial
emotion recognition, presents a particularly challenging recognition
environment owing to its low resolution 48 by 48 pixel grayscale images
collected under diverse and largely uncontrolled real world conditions.
Prior work has demonstrated that achieving accuracy above 70\% is non
trivial, with many baseline CNN architectures plateauing in the 60 to 65\%
range and more sophisticated architectures with extensive hyperparameter
tuning reaching approximately 73\% \cite{b3}. These constraints make
architectural choice and training strategy especially consequential.

This paper presents MoodSync, a deep learning system designed to address
both the accuracy and deployability challenges of real time facial emotion
recognition. MoodSync evaluates three CNN architectures on a balanced five
class subset of FER-2013: a lightweight custom CNN, DenseNet121, and
FERNET, a deep VGG style network. The best performing model, FERNET, is
subsequently deployed in a real time inference pipeline using an OpenCV DNN
based face detector and a temporal prediction smoothing buffer. The system
incorporates a music recommendation module that maps detected emotions to
genre appropriate suggestions, providing a complete affective computing
application.

% -------------------------------------------------------
\section{Literature Review}

Automated facial emotion recognition has become a central challenge in
affective computing and human computer interaction. MoodSync builds upon
prior work in deep learning based FER, spanning traditional CNN pipelines,
attention mechanisms, and hybrid architectures.

\subsection{CNN Based Feature Augmentation}
Amal et al.\ \cite{b1} proposed combining raw pixel inputs with handcrafted
spatial descriptors, specifically face landmarks and Histogram of Oriented
Gradients (HOG) features, before feeding data into a CNN classifier. Face
detection used a Local Binary Pattern (LBP) classifier followed by landmark
localization. The augmented representation yielded 75.1\% accuracy across
seven emotion categories, compared to only 59.1\% with raw pixels alone,
demonstrating that handcrafted descriptors meaningfully complement learned
CNN representations under challenging real world conditions.

\subsection{Spatial Attention Mechanisms}
Abdolrashidi \cite{b2} proposed an attention augmented deep CNN with a
spatial attention module to selectively emphasize discriminative facial
regions such as the eyes, eyebrows, and mouth. Visualization experiments
confirmed that the attention maps aligned with theoretically motivated
facial landmarks. Despite employing fewer layers than many competing
architectures, this attentional model outperformed several standard CNN
baselines on established FER benchmarks.

\subsection{VGGNet Architecture Optimization}
Khaireddin and Chen \cite{b3} demonstrated that a well optimized VGGNet
based classifier, subjected to systematic hyperparameter search, learning
rate scheduling, and optimizer comparison, achieved 73.28\% accuracy on
standard FER datasets. Saliency map analyses confirmed that learned
representations concentrated on the eye and mouth regions. This work directly
motivates MoodSync's adoption of FERNET, a VGG style backbone.

\subsection{Multiscale and Wavelet Channel Attention}
Liu et al.\ \cite{b4} introduced a CNN incorporating multiscale convolutional
layers paired with a wavelet channel attention module to enhance feature
extraction. Multiscale convolutions capture both coarse structural cues and
fine grained textural patterns, while the wavelet based attention module
amplifies informative feature maps across layers. Experiments on two
expression datasets demonstrated consistent accuracy gains over vanilla
CNN and ResNet baselines \cite{b4}.

\subsection{Temporal Modeling via CNN RNN Hybridization}
Manalu and Rifai \cite{b5} developed a hybrid architecture combining CNN
spatial feature extraction with Long Short Term Memory (LSTM) units for
video based emotion detection. Evaluated on the Emotion Wearable dataset,
results confirmed that integrating temporal analysis with spatial feature
extraction substantially broadens the scope of emotion recognition. Although
MoodSync processes individual frames, these findings identify sequential
modeling as a well motivated future direction.

\subsection{Hybrid Deep Classical Classification Pipelines}
Shahzad et al.\ \cite{b6} extracted features from multiple fully connected
layers of pre-trained AlexNet and VGG-16 and supplied them to SVM, KNN, and
ensemble classifiers, achieving accuracy improvements of 7 to 9\% over
standard Softmax baselines on the MLF-W-FER dataset. This resonates with
MoodSync's multi model comparative evaluation strategy.

% -------------------------------------------------------
\section{Methodology}

\subsection{Dataset Description}

MoodSync uses the FER-2013 dataset comprising 48 by 48 pixel grayscale
images collected under diverse real world conditions. This study covers five
emotion categories: angry, happy, neutral, sad, and surprise. The disgust
and fear categories were excluded as they were absent from the local dataset
directory. The full dataset contains 24,176 training images and 6,043 test
images, as detailed in Table~\ref{tab:distribution}.

\begin{table}[htbp]
\caption{Original FER-2013 Class Distribution}
\label{tab:distribution}
\begin{center}
\begin{tabular}{|l|c|c|}
\hline
\textbf{Emotion} & \textbf{Train} & \textbf{Test} \\
\hline
Angry    & 3,995  & 958   \\
Happy    & 7,215  & 1,774 \\
Neutral  & 4,965  & 1,233 \\
Sad      & 4,830  & 1,247 \\
Surprise & 3,171  & 831   \\
\hline
\textbf{Total} & \textbf{24,176} & \textbf{6,043} \\
\hline
\end{tabular}
\end{center}
\end{table}

The training set exhibits a notable class imbalance, with the happy class
containing more than twice the number of images as surprise. Fig.~\ref{fig:samples}
shows one sample image from each emotion class, highlighting the low
resolution and real world diversity of the dataset. Fig.~\ref{fig:dist}
visualizes the full class distribution across train and test splits.

\begin{figure}[htbp]
\centerline{\includegraphics[width=\columnwidth]{emotion_samples.png}}
\caption{Sample images from each of the five emotion classes in the
FER-2013 training set. All images are 48$\times$48 pixels in grayscale.}
\label{fig:samples}
\end{figure}

\begin{figure}[htbp]
\centerline{\includegraphics[width=\columnwidth]{dataset_distribution.png}}
\caption{Dataset distribution: image counts per emotion class for train
and test splits, illustrating class imbalance in the training set.}
\label{fig:dist}
\end{figure}

\subsection{Data Preprocessing}

The preprocessing pipeline applies the following sequential steps
consistently across all three model architectures.

\textit{a) Grayscale Loading and Resizing:} All images were loaded in grayscale
format using OpenCV's \texttt{IMREAD\_GRAYSCALE} flag. A random sample of
100 images per class confirmed uniform dimensions of (48, 48) and a pixel
range of [0, 255]. An explicit resize step was applied to guarantee
consistency across all inputs.

\textit{b) Class Balancing:} The original training set's class imbalance was
addressed by randomly sampling a maximum of 2,000 images per class, resulting
in a balanced training corpus of exactly 10,000 images (2,000 per class). The
test set was used in its original unmodified form (6,043 images) to ensure
evaluation fidelity and avoid introducing optimistic bias.

\textit{c) Normalization:} Pixel values were normalized from the integer range
[0, 255] to floating point values in [0, 1] by dividing by 255.0, which
accelerates gradient descent convergence and ensures numerical stability
during training.

\textit{d) Channel Dimension Expansion:} Normalized arrays were reshaped to
(N, 48, 48, 1) for grayscale models. For DenseNet121, which requires three
channel input, the single channel was replicated three times using NumPy's
\texttt{repeat}, producing tensors of shape (N, 48, 48, 3).

\textit{e) Label Encoding:} String labels were encoded to integer indices using
scikit-learn's \texttt{LabelEncoder} with alphabetical mapping
(angry=0, happy=1, neutral=2, sad=3, surprise=4), then converted to one
hot categorical format using Keras's \texttt{to\_categorical}, yielding
label arrays of shape (N, 5).

\textit{f) Data Augmentation:} Online data augmentation was applied during
training using Keras's \texttt{ImageDataGenerator} with random rotation
($\pm$10 to $\pm$15 degrees), random width and height shift (up to 10\%),
and random horizontal flipping. Augmentation was applied only to the training
pipeline and disabled during test set evaluation. Fig.~\ref{fig:aug} shows
a representative batch of augmented training images.

\begin{figure}[htbp]
\centerline{\includegraphics[width=\columnwidth]{augmented_samples.png}}
\caption{Sample augmented training images illustrating random rotation,
shift, and horizontal flip transformations applied during training.}
\label{fig:aug}
\end{figure}

The final verified input shapes were: X\_train: (10000, 48, 48, 1),
y\_train: (10000, 5), X\_test: (6043, 48, 48, 1), y\_test: (6043, 5).

\subsection{Model Architectures}

\textit{a) FERNET (VGG Style Architecture):} FERNET is the primary proposed
model, inspired by the VGGNet design principle of stacking paired convolutional
layers before spatial downsampling. The architecture comprises three VGG style
blocks followed by a fully connected classifier. Each block contains two Conv2D
layers with Exponential Linear Unit (ELU) activation and Batch Normalization,
followed by Max Pooling and Dropout. Filter depths progress from 64 to 128 to
256 across the three blocks, enabling progressively abstract feature
representations. The classifier consists of Dense(128, ELU), Batch
Normalization, Dropout(0.5), and five class Softmax output. Total trainable
parameters: approximately 2.33 million. Training used the Adam optimizer with
an initial learning rate of 0.001, batch size 128, and a maximum of 50 epochs.
\texttt{EarlyStopping} with patience 12 monitored validation accuracy.
\texttt{ReduceLROnPlateau} halved the learning rate upon validation loss
plateau with patience 5, progressively reducing the rate from 0.001 to
approximately $6.25 \times 10^{-5}$ over the course of training.

\textit{b) Custom CNN (Lightweight Baseline):} The custom CNN is a three
block lightweight baseline. Each block contains one Conv2D with ReLU
activation, Batch Normalization, Max Pooling, and Dropout (0.25). Filter
counts follow 32, 64, and 128. The classifier uses Dense(128, ReLU), Batch
Normalization, Dropout(0.5), and Softmax output. Total parameters:
approximately 357,000. Training used batch size 64, up to 60 epochs, and
\texttt{EarlyStopping} with patience 5. Early stopping triggered at Epoch 15.

\textit{c) DenseNet121 (Deep Architecture Baseline):} DenseNet121 was
evaluated without pretrained ImageNet weights and without its top
classification layers. A \texttt{GlobalAveragePooling2D} layer was appended,
followed by Dense(256, ELU), Batch Normalization, Dropout(0.5), and Softmax
output. Total parameters: approximately 7.30 million. Training used batch
size 64, up to 40 epochs, \texttt{EarlyStopping} with patience 7, and
\texttt{ReduceLROnPlateau} with patience 3. All grayscale inputs were
converted to pseudo-RGB by channel replication. Early stopping triggered at
Epoch 27.

\subsection{Evaluation Metrics}

All models were evaluated on the unmodified FER-2013 test set using overall
accuracy, weighted precision, weighted recall, weighted F1 score, and macro
average AUC. Weighted metrics account for test set class imbalance by
weighting per class scores by the number of true instances in each class.
Macro average AUC was computed by averaging the one vs.\ rest AUC values
across all five emotion classes using scikit-learn's \texttt{roc\_curve}
and \texttt{auc} functions.

\subsection{Real Time Inference Pipeline}

Following model evaluation, FERNET was deployed in a real time webcam
inference pipeline. Face detection used OpenCV's DNN module with a ResNet-10
SSD based face detector. Face regions with detection confidence below 0.5
were discarded. Each accepted face crop was converted to grayscale, resized
to 48 by 48 pixels, normalized, and reshaped to (1, 48, 48, 1) for
prediction. A sliding window prediction smoothing buffer of 15 frames was
maintained using Python's \texttt{collections.deque}. The final emotion
label at each frame was determined by averaging the predicted probability
vectors across the buffer and selecting the class with the highest mean
probability. Each detected emotion was then mapped to a genre appropriate
music recommendation via a fixed lookup table.

% -------------------------------------------------------
\section{Results and Discussion}

\subsection{Performance Metrics}

Table~\ref{tab:results} presents the complete performance metrics for all
three models on the FER-2013 test set. FERNET achieved the highest scores
across all five metrics. Fig.~\ref{fig:perf} visualizes the accuracy and
F1 score comparison, showing a gap of approximately 6.65 percentage points
between FERNET and the Custom CNN, and 15.43 percentage points between
FERNET and DenseNet121 in accuracy.

\begin{table}[htbp]
\caption{Model Performance Comparison on FER-2013 Test Set}
\label{tab:results}
\begin{center}
\begin{tabular}{|l|c|c|c|c|c|}
\hline
\textbf{Model} & \textbf{Acc.} & \textbf{Prec.} & \textbf{Rec.} & \textbf{F1} & \textbf{AUC} \\
\hline
FERNET      & 71.45\% & 72.28\% & 71.45\% & 71.60\% & 0.922 \\
Custom CNN  & 64.80\% & 65.22\% & 64.80\% & 64.80\% & 0.887 \\
DenseNet121 & 56.02\% & 58.93\% & 56.02\% & 56.15\% & 0.848 \\
\hline
\end{tabular}
\end{center}
\end{table}

\begin{figure}[htbp]
\centerline{\includegraphics[width=\columnwidth]{model_performance.png}}
\caption{Final model performance comparison: accuracy and F1 score for
FERNET, Custom CNN, and DenseNet121 on the FER-2013 test set.}
\label{fig:perf}
\end{figure}

\subsection{Model Analysis}

FERNET trained over the full 50 epochs with consistent improvement in
validation accuracy. The progressive learning rate reduction strategy proved
effective: each scheduled halving of the learning rate produced a measurable
accuracy gain. Training accuracy reached 77.18\% at Epoch 50 against a
validation accuracy of 71.45\%, indicating moderate overfitting controlled by
Dropout regularization and online data augmentation. The high AUC of 0.922
confirms strong discriminative capability across all five emotion classes,
consistent with Khaireddin and Chen \cite{b3} who achieved 73.28\% on a
comparable benchmark with a similarly optimized VGGNet variant.

The Custom CNN converged rapidly, with early stopping triggered at Epoch 15.
Despite its substantially smaller parameter count (357K versus 2.33M for
FERNET), it achieved a solid 64.80\% accuracy. The alignment between accuracy
(64.80\%) and F1 score (64.80\%) indicates balanced performance across emotion
classes without systematic bias.

DenseNet121 achieved the lowest accuracy at 56.02\% despite having the
highest parameter count (7.30 million). This underperformance is attributable
to two principal factors. First, DenseNet121 was designed for higher
resolution inputs, and its dense connectivity pattern adapts less effectively
to 48 by 48 pixel FER-2013 images. Second, training from random
initialization required learning 7.30 million parameters from only 10,000
training images, which is insufficient for a network of this depth.
Significant validation loss spikes at Epoch 3 (29.08) and Epoch 5 (38.49)
indicate training instability from which the model only partially recovered.

\subsection{Confusion Matrix Analysis}

Fig.~\ref{fig:cm1}, Fig.~\ref{fig:cm2}, and Fig.~\ref{fig:cm3} show the
confusion matrices for FERNET, Custom CNN, and DenseNet121 respectively.
Several consistent patterns are evident across all three models.

\begin{figure}[htbp]
\centerline{\includegraphics[width=\columnwidth]{cm_fernet.png}}
\caption{Confusion matrix for FERNET on the FER-2013 test set.}
\label{fig:cm1}
\end{figure}

The happy class is the best classified emotion in all models. FERNET
correctly classified 1,453 of 1,774 happy samples (81.9\%), as the happy
expression has highly distinctive visual features including a prominent smile
and raised cheeks. The surprise class also achieved strong results in FERNET,
with 712 of 831 samples correctly classified (85.7\%), the highest per class
recall for any emotion in FERNET.

\begin{figure}[htbp]
\centerline{\includegraphics[width=\columnwidth]{cm_custom_cnn.png}}
\caption{Confusion matrix for Custom CNN on the FER-2013 test set.}
\label{fig:cm2}
\end{figure}

The sad class exhibits the greatest misclassification rate across all models.
In FERNET, only 676 of 1,247 sad samples (54.2\%) are correctly classified,
with 315 misclassified as neutral, the largest single off diagonal entry in
the FERNET matrix. This overlap between sad and neutral is consistent with
the literature, as both expressions involve similar reduced facial muscle
activation \cite{b2}. The neutral class similarly suffers from confusion with
sad, with FERNET correctly classifying 877 of 1,233 neutral samples (71.1\%).

\begin{figure}[htbp]
\centerline{\includegraphics[width=\columnwidth]{cm_densenet121.png}}
\caption{Confusion matrix for DenseNet121 on the FER-2013 test set.}
\label{fig:cm3}
\end{figure}

The angry class shows moderate performance in FERNET (600 of 958 correct,
62.6\%), with frequent confusion with both neutral (149) and sad (146),
reflecting the subtle visual differences between these states in low
resolution grayscale images.

\subsection{ROC Curve Analysis}

Fig.~\ref{fig:roc} presents the macro average ROC curves for all three
models. All models perform well above the random classifier baseline (the
dashed diagonal line), demonstrating genuine discriminative capability.
FERNET's curve (blue) rises steeply from the origin, achieving a true
positive rate of approximately 0.78 at a false positive rate of only 0.10,
confirming high sensitivity while minimizing false alarms across all emotion
classes. The Custom CNN (green) achieves a slightly lower but competitive AUC
of 0.887, while DenseNet121 (red) trails with 0.848. The ordering of the
three curves is consistent with the accuracy and F1 rankings in
Table~\ref{tab:results}.

\begin{figure}[htbp]
\centerline{\includegraphics[width=\columnwidth]{roc_curves.png}}
\caption{Macro average ROC curves for all three models. FERNET achieves
the highest AUC of 0.922, followed by Custom CNN (0.887) and DenseNet121
(0.848).}
\label{fig:roc}
\end{figure}

\subsection{Real Time System and Music Recommendation}

The deployed FERNET pipeline, integrated with the ResNet-10 SSD face
detector and 15 frame temporal smoothing buffer, demonstrated stable and
responsive emotion classification under real time webcam conditions. The
smoothing buffer effectively reduced frame to frame label flickering,
producing a consistent and natural user experience. Table~\ref{tab:music}
summarizes the emotion to music genre mapping used in deployment.

\begin{table}[htbp]
\caption{Emotion to Music Recommendation Mapping}
\label{tab:music}
\begin{center}
\begin{tabular}{|l|l|}
\hline
\textbf{Detected Emotion} & \textbf{Music Recommendation} \\
\hline
Angry    & Relaxing Music --- Lo-fi / Acoustic \\
Happy    & Upbeat Music --- Pop / Dance / EDM  \\
Neutral  & Ambient Playlist --- Chill / Lounge \\
Sad      & Calm Music --- Piano / Soft Acoustic \\
Surprise & Energetic Music --- Rock / Hype Beats \\
\hline
\end{tabular}
\end{center}
\end{table}

To validate the end to end functionality of the MoodSync system, the
deployed pipeline was tested in a live two person webcam session. Three
consecutive emotional states were enacted and captured: neutral, surprise,
and happy. In each case the system correctly detected both faces using the
ResNet-10 SSD detector, assigned emotion labels with associated confidence
scores, and displayed the corresponding music recommendation on screen.

Fig.~\ref{fig:rt_neutral} shows the neutral detection state. Both subjects
exhibit a composed, expressionless appearance. The system correctly assigned
the neutral label to both detected faces with a smoothed confidence of 75.8\%,
triggering the Ambient Playlist recommendation.

\begin{figure}[htbp]
\centerline{\includegraphics[width=\columnwidth]{realtime_neutral.png}}
\caption{Real time inference output for the neutral emotion. Both faces are
correctly detected and labelled as Neutral (75.8\%) with the Ambient
Playlist music recommendation displayed.}
\label{fig:rt_neutral}
\end{figure}

Fig.~\ref{fig:rt_surprise} shows the surprise detection state. Both subjects
exhibited wide eyes and open mouths. The system assigned the surprise label
to both faces with confidence scores of 58.6\% and 59\% respectively, and
triggered the Energetic Music recommendation. The lower confidence compared
to the neutral state reflects the inherent ambiguity of the surprise
expression at real world webcam resolution.

\begin{figure}[htbp]
\centerline{\includegraphics[width=\columnwidth]{realtime_surprise.png}}
\caption{Real time inference output for the surprise emotion. Both faces are
correctly labelled as Surprise with confidence scores of 58.6\% and 59\%,
and the Energetic Music recommendation is displayed.}
\label{fig:rt_surprise}
\end{figure}

Fig.~\ref{fig:rt_happy} shows the happy detection state. Both subjects
displayed prominent smiles. The system achieved the highest per frame
confidences across all three demonstration states, correctly classifying
both faces as happy with scores of 86.9\% and 89.6\% respectively, and
displaying the Upbeat Music recommendation. This result is consistent with
the confusion matrix findings in Section~IV-C, which identified happy as
the best classified emotion across all evaluated models.

\begin{figure}[htbp]
\centerline{\includegraphics[width=\columnwidth]{realtime_happy.png}}
\caption{Real time inference output for the happy emotion. Both faces are
correctly labelled as Happy with high confidence scores of 86.9\% and
89.6\%, and the Upbeat Music recommendation is displayed.}
\label{fig:rt_happy}
\end{figure}

Across all three demonstrations, the system accurately detected multiple
faces simultaneously, applied temporal smoothing to suppress frame level
flickering, and rendered the correct music recommendation in real time.
The colour coded bounding box scheme, where cyan indicates neutral, yellow
indicates surprise, and green indicates happy, provides an immediate visual
indicator of the detected emotional state, enhancing usability.

% -------------------------------------------------------
\section{Conclusion}

This paper presented MoodSync, a real time facial emotion recognition system
integrating a deep VGG style CNN with an adaptive music recommendation
module. Three CNN architectures were systematically evaluated on a balanced
subset of the FER-2013 dataset. FERNET achieved the best performance at
71.45\% accuracy, 71.60\% weighted F1 score, and AUC of 0.922, substantially
outperforming the Custom CNN (64.80\%) and DenseNet121 (56.02\%). Confusion
matrix analysis revealed that happy and surprise are the best classified
emotions, while sad and neutral exhibit the highest misclassification overlap
due to their visual similarity at low resolution. ROC analysis confirmed that
all three models operate well above the random baseline, with FERNET
demonstrating the strongest overall discriminability. The real time deployment
pipeline, combining DNN based face detection with a 15 frame prediction
smoothing buffer, demonstrated stable and responsive performance, confirming
MoodSync's practical viability as an end to end affective computing application.

Several directions for future work were identified. First, training used only
10,000 of the available 24,176 training images; incorporating all data with
weighted loss functions or SMOTE oversampling may yield further accuracy
improvements. Second, integrating recurrent structures such as LSTM \cite{b5}
could improve temporal consistency for continuous video streams. Third, the
music recommendation module could be replaced with a personalized engine
accounting for user preferences and listening history. Finally, extending
the emotion taxonomy to include disgust and fear would broaden coverage for
general affective computing applications.

% -------------------------------------------------------
\begin{thebibliography}{00}

\bibitem{b1}
V.~S. Amal, S.~Suresh, and G.~Deepa,
``Real-time emotion recognition from facial expressions using convolutional
neural network with FER2013 dataset,''
in \textit{Smart Innovation, Systems and Technologies},
pp. 541--551, Springer Singapore, 2021.

\bibitem{b2}
S.~M. M. M.~A. Abdolrashidi,
``Deep-Emotion: Facial expression recognition using attentional convolutional
network,''
\textit{Sensors}, vol. 21, no. 9, p. 3046, 2021.

\bibitem{b3}
Y.~Khaireddin and Z.~Chen,
``Facial emotion recognition: State of the art performance on FER2013,''
\textit{arXiv preprint arXiv:2105.03588}, 2021.

\bibitem{b4}
J.~Liu, X.~Lin, P.~Ji, J.~Chen, and J.~Zhang,
``Multiscale wavelet attention convolutional network for facial expression
recognition,''
\textit{Scientific Reports}, vol. 15, no. 1, p. 22219, 2025.

\bibitem{b5}
H.~V. Manalu and A.~P. Rifai,
``Detection of human emotions through facial expressions using hybrid
convolutional neural network recurrent neural network algorithm,''
\textit{Intelligent Systems with Applications}, vol. 21, p. 200339, 2024.

\bibitem{b6}
H.~M. Shahzad, S.~M. Bhatti, A.~Jaffar, S.~Akram, M.~Alhajlah,
and A.~Mahmood,
``Hybrid facial emotion recognition using CNN-based features,''
\textit{Applied Sciences}, vol. 13, no. 9, p. 5572, 2023.

\end{thebibliography}

\end{document}
